\subsection{Summary}

This thesis set out to build skin-lesion classifiers for ISIC-2019 that are both
accurate across all eight classes and practical to deploy, and that use the
clinical metadata a dermatologist would normally have on hand. We approached the
problem with a single design philosophy: convolutional features for local
texture, a transformer trunk for global structure, and a cross-attention head
that folds in age, sex, and lesion site, realised at two points on the
efficiency curve. The lightweight hybrid, at 3.98M parameters and 0.631 GMACs,
reaches a balanced multi-class accuracy of 0.6081 and is the most
parameter-efficient model in the sub-6M budget, beating all five lightweight
baselines including the strong hybrid baselines MobileViT-S and
EfficientFormer-L1. The flagship DEKAN model, at 16.45M parameters, reaches 0.6438 BMA, comparable
to the balanced accuracy of the top ISIC-2019 leaderboard ensemble in a single
model; a linear-head ablation of the same architecture (dekan\_linear) reaches
0.6510, the highest figure in the study, with the implication that the KAN
classifier head adds no accuracy benefit on this task.

\subsection{Future Work}

The most immediate next steps follow directly from the limitations in
Chapter~5. Repeating every run across at least three seeds and reporting mean and
standard deviation would put the comparisons on a publication-grade footing.
Completing the two remaining unfinished ablations (DenseNet-121 alone and
EfficientNet-B0 alone within the DEKAN framework) would pin down the
individual contribution of each backbone and quantify the size of the
dual-backbone fusion gain. Evaluating on an external dataset, for example mapping
the eight-class output to a melanoma-versus-rest decision on a separate
collection, would test generalisation beyond ISIC-2019. Finally, qualitative
interpretability, using Grad-CAM on the CNN stem and attention rollout on the
transformer trunk, would help confirm that the models attend to clinically
meaningful regions, which matters for any tool intended to support, rather than
replace, a clinician.
